\documentclass[11pt]{letter}
\usepackage[utf8]{inputenc}
\usepackage{geometry}
\geometry{letterpaper, margin=1in}
\usepackage{hyperref}
\usepackage{graphicx}

\signature{Nam Nguyen \\ Principal Investigator, Lead AI Researcher}
\address{
    Department of Biomedical Engineering \\
    [Institution Name] \\
    [City, State, Zip] \\
    [Email Address]
}

\begin{document}

\begin{letter}{
    Prof. [Editor-in-Chief Name] \\
    Editor-in-Chief \\
    IEEE Transactions on Medical Imaging (T-MI)
}

\opening{Dear Prof. [Editor-in-Chief Name],}

We are pleased to submit our original research manuscript entitled, \textbf{"Breaking the Feedback Trap: Learned Residual Decoupling and the Pareto Frontier in Medical Image Segmentation"} for consideration for publication in \textit{IEEE Transactions on Medical Imaging}.

In this work, we present a fundamental investigation into a pervasive structural pathology in recurrent medical image segmentation networks, which we define as the \textit{Feedback Trap}. We demonstrate that standard iterative spatial guidance inevitably leaks gradients backward through recurrent history, causing catastrophic error accumulation and severe false-positive over-segmentation. 

Our core methodological contribution is the introduction of \textbf{Learned Residual Decoupling}, a gating mechanism that strictly severs backward gradient propagation while restoring complete architectural controllability via a $1\times1$ Convolution Gate. This mechanism operates with zero parameter or FLOP overhead compared to standard models, yet stabilizes iterative learning.

Crucially, rather than superficially claiming artificial state-of-the-art metrics, our manuscript rigorously dissects the optimization landscape. By demonstrating how our network perfectly obeys bidirectional loss mandates (aggressively pruning FPR in Phase 7C or recovering Recall in Phase 7D), we expose the inherent \textit{Precision-Recall Pareto Frontier} of the challenging Kvasir-SEG Sessile dataset. We prove both mathematically and empirically that standard CNN intrinsic features have hit a physical boundary on flat polyps due to extreme boundary ambiguity, paving the way for future temporal sequence or multimodal Foundation Model research.

We believe that our theoretical diagnostics, robust empirical evaluations (including zero-shot cross-domain transfers), and deep intellectual honesty regarding dataset limitations align strongly with the rigorous standards of T-MI. 

This manuscript has not been published and is not under consideration for publication elsewhere. All authors have read and approved the manuscript and take full responsibility for its content.

Thank you for your time and consideration.

\closing{Sincerely,}

\end{letter}
\end{document}
